(1) 静水压关系 $dp/dz=-\rho g$，理想气体 $p=\rho R_sT$。每单位质量的可逆绝热第一定律可写 $c_p dT=(1/\rho)dp$，代入压强梯度得 $c_p\,dT/dz=-g$，所以干绝热直减率 $\Gamma_d=-dT/dz=g/c_p$，与 $R_s$ 是否取摩尔量无关，只需一致使用比热单位。
(2) 积分得 $T(z)=T_0-gz/c_p$。理想气体与绝热式联立 $d\ln p=(c_p/R_s)\,d\ln T$，于是 $p(z)=p_0[T(z)/T_0]^{c_p/R_s}$；仅 $T(z)>0$ 且气体性质、重力常数近似有效的高度区可使用此幂律，绝不能延拓到数学上负温度。
(3) 对小气团的绝热变化率 $dT_{\rm parcel}/dz=-g/c_p$，而背景变化为 $T_{\rm env}'(z)$。向上微移 $\xi$ 后的密度差在相同压强处约为 $\delta\rho/\rho\simeq-\delta T/T$，且 $\delta T=[-g/c_p-T_{\rm env}']\xi$。因而恢复加速度 $\ddot\xi=-N^2\xi$，其中 $N^2=(g/T_{\rm env})(T_{\rm env}'+g/c_p)$。
(4) 等温背景 $T_{\rm env}'=0$ 时 $N^2=g^2/(c_pT)>0$，小位移振荡稳定；干绝热背景 $T_{\rm env}'=-g/c_p$ 时 $N^2=0$，理想中性浮力稳定；若温度随高度下降得更快，即 $T_{\rm env}'<-g/c_p$，则 $N^2<0$，微扰指数增长引发对流不稳定。不能把气团自身的绝热直减率直接当作任意环境的实际直减率。
